\scanpage{213}

où $\underline{\mathrm{Cart}}_{X\unsure{,|/}S}$ est le faisceau de Cartier déjà envisagé dans (3.2). Regardons le sous foncteur de \deleted{$\unsure{\mathrm{Pic}}$} $P^{cin}_{X/S}$ image inverse de $\underline{\mathrm{Pic}}^0_{X/S}$ et supposons que $\underline{\mathrm{Pic}}^0_{X/S}$ est $V$-divisible (p.ex. un schéma abélien), de sorte que $(\underline{\mathrm{Pic}}^{0\deleted{\unsure{p/S}}}_{X/S})^{p/S}\unsure{\longleftarrow|\longrightarrow|\text{—}}\underline{\mathrm{Pic}}^0_{X/S}$ est un épimorphisme. On trouve donc une suite exacte

\[
0\longrightarrow U\longrightarrow [P^{cin}_{X/S}]^{0\ \substack{\text{(image inverse}\\\text{de }\underline{\mathrm{Pic}}^0_{X/S})}}\longrightarrow\underline{\mathrm{Pic}}^0_{X/S}\longrightarrow0\tag*{(4.6)}
\]

\uline{NB} Une façon commode d'exprimer $P^{cin}_{X/S}$ est de dire que c'est le complété formel de $P^{cip}_{X/S}\simeq\underline{\mathrm{Pic}}_{X/S}^{\deleted{\unsure{p}}(p/S)}\subset P^{ci}$ dans $P^{ci}$. Dans le cas favorable où $P^{ci\,\unsure{0|o}}$ est l'extension vectorielle universelle \inserted{$E$} d'un schéma abélien $\unsure{B|\mathscr B}$ par un grpe vectoriel $V$, on trouve \unsure{par|pour|que} \deleted{$P^{cin\,0}$} donc $B^{p/S}\subset E$, et on \deleted{trouve} \inserted{pour $P^{cin\,0}_{X/S}$} le complété formel \inserted{le long} \unsure{dudit|du dit|déduit|induit} $B^{p/S}$. Sa connaissance est donc nettement plus précise que celle du complété formel \inserted{$\overline E$} de $E$ le long de sa section unité, i.e. le groupe formel associé à $E$.

% archive sheet number: 212
